\documentclass[11pt,a4paper]{article}

% ------------------------------------------------------------
% Grundlayout
% ------------------------------------------------------------
\usepackage{fontspec}
\setmainfont{Latin Modern Roman}
% Die deutsche Sprachinstallation von Babel ist auf manchen TinyTeX-Systemen
% nicht als Paketoption registriert; die benötigten Überschriften werden daher
% direkt gesetzt.
\renewcommand{\contentsname}{Inhaltsverzeichnis}
\renewcommand{\figurename}{Abbildung}
\renewcommand{\tablename}{Tabelle}
\usepackage[a4paper,top=2.5cm,bottom=2.5cm,left=2.7cm,right=2.7cm]{geometry}
\usepackage{xcolor}
\definecolor{ADAblue}{HTML}{173F5F}
\definecolor{ADAgold}{HTML}{D99A2B}
\definecolor{ADAlight}{HTML}{EEF3F7}
\usepackage{booktabs}
\usepackage{longtable}
\usepackage{array}
\usepackage{tabularx}
\usepackage{graphicx}
\usepackage{hyperref}
\hypersetup{
  colorlinks=true,
  linkcolor=ADAblue,
  urlcolor=ADAblue,
  citecolor=ADAblue,
  pdftitle={E-Portfolio Santander Customer Transaction Prediction},
  pdfauthor={Cyprian Sawarzynski}
}

% ------------------------------------------------------------
% Eigene Formatierung und Platzhalter
% ------------------------------------------------------------
\newcommand{\placeholder}[1]{\textcolor{ADAgold}{\textbf{[PLACEHOLDER: #1]}}}
\newcommand{\todo}[1]{\textcolor{red}{\textbf{[TODO: #1]}}}
\newcommand{\logentry}[2]{\subsubsection*{#1 – #2}}
\newcommand{\sourceitem}[2]{\noindent\textbf{#1}\\#2\par\medskip}

\pagestyle{plain}

\begin{document}

% ------------------------------------------------------------
% Titelseite
% ------------------------------------------------------------
\begin{titlepage}
  \centering
  \vspace*{2.2cm}
  {\color{ADAblue}\rule{\textwidth}{1.5pt}}\\[1.2cm]
  {\Huge\bfseries\color{ADAblue} E-Portfolio\par}
  \vspace{0.5cm}
  {\LARGE Santander Customer Transaction Prediction\par}
  \vspace{1.4cm}
  {\color{ADAgold}\rule{0.35\textwidth}{2pt}}\\[1.4cm]

  \begin{tabular}{>{\bfseries}p{0.34\textwidth}p{0.48\textwidth}}
    Name & Cyprian Sawarzynski \\
    Modul & Advanced Data Analytics / Scientific Programming and Data Analysis \\
    Semester & Sommersemester 2026 \\
    Gruppennummer & 138 \\
    Gruppenmitglieder & Tim Hebestreit (Gruppenleiter), Angeline Usanayo, Trong Cao \\
    Sprecher/Gruppenleitung & Tim Hebestreit \\
    Abgabetermin & 01.09.2026 \\
  \end{tabular}

  \vfill
  {\small Individuelles E-Portfolio · Version: \today\par}
\end{titlepage}

\tableofcontents
\clearpage

% ------------------------------------------------------------
% 1 Projektüberblick
% ------------------------------------------------------------
\section{Projektüberblick und Fragestellung}

\subsection{Datensatz und Aufgabenstellung}
Das Projekt behandelt den Datensatz \emph{Santander Customer Transaction Prediction}. Die
Daten stammen aus einer Kaggle Competition aus dem Jahr 2019. Auf Basis von
Finanztransaktionsmerkmalen soll vorhergesagt werden, ob eine Person eine Transaktion
durchführen wird.

\begin{itemize}
  \item OpenML-Datensatz: ID \texttt{45566}
  \item Kaggle: \url{https://www.kaggle.com/c/santander-customer-transaction-prediction}
  \item OpenML: \url{https://www.openml.org/search?type=data\&status=active\&id=45566}
\end{itemize}

\subsection{Eigene Forschungsfrage}
\placeholder{präzise eigene Forschungsfrage ergänzen, z.~B. welche Merkmale und Modelle
die Transaktionsvorhersage besonders gut unterstützen.}

\subsection{Projektziel}
Ziel ist die Entwicklung, der Vergleich und die begründete Evaluation mehrerer
Machine-Learning-Modelle für die binäre Klassifikationsaufgabe. Die Entscheidungen sollen
durch die explorative Datenanalyse, geeignete Daten-Pipelines, passende Metriken und
dokumentierte Modellvergleiche nachvollziehbar begründet werden.

% ------------------------------------------------------------
% 2 Logbuch
% ------------------------------------------------------------
\section{Individuelles Logbuch}
Das vollständige, fortlaufend gepflegte Logbuch ist im Anhang eingebunden. Die separate
Quelldatei bleibt dabei die einzige Stelle, an der die Einträge gepflegt werden.

% ------------------------------------------------------------
% 3 Zusammenfassung
% ------------------------------------------------------------
\section{Zusammenfassung des Logbuchs}
\textit{Dieser Abschnitt wird am Ende aus dem Logbuch entwickelt und auf maximal vier
Seiten einschließlich Abbildungen und Tabellen verdichtet.}

\subsection{Fragestellung und Ziel}
\placeholder{Fragestellung, fachliche Motivation und Ziel der Analyse ausformulieren.}

\subsection{Vorgehen und Methodeneinsatz}
\placeholder{Datenzugang, Datentypen, EDA, Umgang mit fehlenden Werten, Train-Test-Aufteilung,
Pipelines, Modelle, Learning Curves, Dimensionsreduktion und Hyperparameteroptimierung
beschreiben.}

\subsection{Iterationen und begründete Entscheidungen}
\placeholder{Mehrere getestete Ansätze gegenüberstellen und anhand der Evaluationsergebnisse
begründen.}

\subsection{Ergebnisse}
\placeholder{Zentrale Kennzahlen, Modellvergleich und wichtigste EDA-Befunde eintragen.}

\begin{figure}[htbp]
  \centering
  \fbox{\parbox[c][4cm][c]{0.82\textwidth}{\centering
    \textcolor{ADAgold}{\textbf{PLACEHOLDER: Abbildung einfügen}}\\
    z.~B. Modellvergleich, ROC-/PR-Kurve oder EDA-Ergebnis}}
  \caption{\placeholder{aussagekräftige Bildunterschrift}}
  \label{fig:placeholder}
\end{figure}

\subsection{Grenzen, Unsicherheiten und Ausblick}
\placeholder{Datenqualität, Rechenzeit, Generalisierbarkeit, mögliche Verzerrungen und
konkrete nächste Schritte reflektieren.}

% ------------------------------------------------------------
% 4 Einschätzung
% ------------------------------------------------------------
\section{Einschätzung der Gruppenmitglieder}
Dieser vertrauliche Abschnitt reflektiert die Zusammenarbeit. Für jede Person werden
mindestens eine Stärke und ein Entwicklungsfeld mit kurzer Begründung ergänzt.

\subsection*{Tim Hebestreit}
\textbf{Stärke:} \placeholder{konkrete Stärke und Begründung}\\
\textbf{Entwicklungsfeld:} \placeholder{konkretes Entwicklungsfeld und Begründung}

\subsection*{Angeline Usanayo}
\textbf{Stärke:} \placeholder{konkrete Stärke und Begründung}\\
\textbf{Entwicklungsfeld:} \placeholder{konkretes Entwicklungsfeld und Begründung}

\subsection*{Trong Cao}
\textbf{Stärke:} \placeholder{konkrete Stärke und Begründung}\\
\textbf{Entwicklungsfeld:} \placeholder{konkretes Entwicklungsfeld und Begründung}

% ------------------------------------------------------------
% 5-6 Gruppenabschnitte
% ------------------------------------------------------------
\section{Ergebnisse der Gruppe}
\textbf{Hinweis:} Dieser Abschnitt ist laut Aufgabenstellung nur im Portfolio der
Sprecherin enthalten. Tim Hebestreit führt als Gruppenleiter die gemeinsame Darstellung.

\placeholder{Falls von Tim freigegeben, hier auf die finale Gruppenzusammenfassung und
zentrale gemeinsame Ergebnisse verweisen.}

\section{Protokoll der Gruppentreffen}
\textbf{Hinweis:} Dieser Abschnitt ist laut Aufgabenstellung nur im Portfolio der
Sprecherin enthalten. Eigene Bezüge zu Treffen werden zusätzlich im Logbuch dokumentiert.

\placeholder{Verweis auf Tims Protokoll bzw. relevante Treffen ergänzen, sofern erforderlich.}

% ------------------------------------------------------------
% 7 Quellen und Hilfsmittel
% ------------------------------------------------------------
\section{Quellen und Hilfsmittel}
Zitierstil: \placeholder{z.~B. APA 7 oder Harvard verbindlich festlegen}

\sourceitem{Waack, J.-M. (2026).}{\emph{Bewertungshinweise ADA}, SoSe 2026. Verwendung:
Anforderungen und Bewertung des E-Portfolios. Zugriff: Projektunterlagen.}

\sourceitem{OpenML, Datensatz ID 45566.}{\emph{Santander Customer Transaction Prediction}.
Verwendung: Datenquelle. Zugriff: \placeholder{Zugriffsdatum ergänzen}.}

\sourceitem{Kaggle.}{\emph{Santander Customer Transaction Prediction}. Verwendung:
Wettbewerbskontext und Datensatzbeschreibung. Zugriff: \placeholder{Zugriffsdatum ergänzen}.}

\sourceitem{Software.}{Python, Jupyter, pandas, NumPy, scikit-learn, matplotlib und seaborn.
Verwendete Versionen: \placeholder{Versionen eintragen}.}

\sourceitem{Generative KI.}{Tool, Zweck und relevante Prompts:
\placeholder{transparent dokumentieren}. Datum: \placeholder{Datum}.}

\subsection{KI- und Tool-Transparenz}
Für jeden relevanten Einsatz werden Datum, Tool/Version, Zweck, Prompt bzw. Kurzbeschreibung,
Ergebnis und die eigene Prüfung des Ergebnisses dokumentiert.

% ------------------------------------------------------------
% Anhang: Checkliste
% ------------------------------------------------------------
\appendix
\section{Individuelles Logbuch}
\def\PortfolioMain{}
\ifdefined\PortfolioMain
\else
\documentclass[11pt,a4paper]{article}

% Eigenständiges individuelles Logbuch für das Santander-Projekt.
\usepackage{fontspec}
\setmainfont{Latin Modern Roman}
\usepackage[a4paper,top=2.5cm,bottom=2.5cm,left=2.7cm,right=2.7cm]{geometry}
\usepackage{xcolor}
\definecolor{ADAblue}{HTML}{173F5F}
\definecolor{ADAgold}{HTML}{D99A2B}
\usepackage{booktabs}
\usepackage{longtable}
\usepackage{array}
\usepackage{tabularx}
\usepackage{hyperref}
\hypersetup{
  colorlinks=true,
  linkcolor=ADAblue,
  urlcolor=ADAblue,
  pdftitle={Logbuch Santander Customer Transaction Prediction},
  pdfauthor={Cyprian Sawarzynski}
}

\newcommand{\placeholder}[1]{\textcolor{ADAgold}{\textbf{[PLACEHOLDER: #1]}}}
\newcommand{\logsection}[2]{\subsection*{#1 -- #2}}
\newcommand{\field}[2]{\noindent\textbf{#1}\\#2\par\medskip}

\title{\color{ADAblue}\textbf{Individuelles Logbuch}\\[0.4em]
       \large Santander Customer Transaction Prediction}
\author{Cyprian Sawarzynski\\Gruppe 138}
\date{Sommersemester 2026}
\fi

\providecommand{\placeholder}[1]{\textcolor{ADAgold}{\textbf{[PLACEHOLDER: #1]}}}
\providecommand{\logsection}[2]{\subsection*{#1 -- #2}}
\providecommand{\field}[2]{\noindent\textbf{#1}\\#2\par\medskip}

\ifdefined\PortfolioMain
\else
\begin{document}
\maketitle
\tableofcontents
\clearpage
\fi

\section{Zweck und Führungsregeln}

Dieses Logbuch dokumentiert meinen individuellen Arbeitsprozess während des gesamten
Bearbeitungszeitraums. Es ist kein nachträglich geglätteter Ergebnisbericht, sondern hält
den tatsächlichen Weg zur Analyse fest. Dadurch sollen Arbeitsaufwand, eigene Entscheidungen,
Anpassungen, Schwierigkeiten und verworfene Ansätze nachvollziehbar bleiben.

Die Einträge werden chronologisch und möglichst unmittelbar nach dem jeweiligen Arbeitsschritt
ergänzt. Jeder Eintrag enthält mindestens:

\begin{itemize}
  \item Datum des Eintrags,
  \item einen kurzen und prägnanten Titel,
  \item eine verständliche Zusammenfassung des Arbeitsschritts.
\end{itemize}

Zusätzlich dokumentiere ich systematisch:

\begin{itemize}
  \item konkrete Arbeitsschritte und Ergebnisse,
  \item methodische Entscheidungen und deren Begründung,
  \item verworfene Ansätze und Iterationen,
  \item Schwierigkeiten, Unsicherheiten und offene Fragen,
  \item Zeitaufwand und gegebenenfalls Bezug zu Gruppentreffen,
  \item nächste Schritte und Abweichungen vom vereinbarten Vorgehen.
\end{itemize}

Das Logbuch bildet die Grundlage für die spätere Zusammenfassung im E-Portfolio. Abbildungen
oder Tabellen dürfen hier auch vorläufig bzw. als Arbeitsdokumentation eingefügt werden;
für die finale Zusammenfassung werden die relevanten Darstellungen später formal ausgearbeitet.

\section{Projekt- und Gruppenrahmen}

\begin{tabularx}{\textwidth}{>{\bfseries}p{0.32\textwidth}X}
  Thema & Santander Customer Transaction Prediction \\
  Datensatz & OpenML ID 45566 \\
  Name & Cyprian Sawarzynski \\
  Gruppennummer & 138 \\
  Gruppenmitglieder & Tim Hebestreit (Gruppenleiter), Angeline Usanayo, Trong Cao \\
  Sprecher/Gruppenleitung & Tim Hebestreit \\
  Projektdeadline & 01.09.2026 \\
  Individueller Workload-Richtwert & ca. 100 Stunden \\
\end{tabularx}

\section{Chronologische Einträge}

\logsection{06.08.2026}{Projektunterlagen und Projektauswahl}
\field{Zusammenfassung}{Die Bewertungs- und Projektunterlagen wurden gelesen. Als Projekt aus
den vorhandenen Aufgaben wurde Santander Customer Transaction Prediction ausgewählt.}
\field{Arbeitsschritte}{
\begin{itemize}
  \item Bewertungshinweise für das E-Portfolio gelesen.
  \item Anforderungen an EDA und Machine Learning herausgearbeitet.
  \item Santander-Aufgabenblatt gelesen und OpenML-Datensatz ID 45566 notiert.
  \item Kaggle- und OpenML-Links als spätere Quellen vorgemerkt.
\end{itemize}}
\field{Entscheidungen und Begründung}{Die Analyse wird auf Santander fokussiert, weil die Aufgabe eine
binäre Transaktionsvorhersage mit klarer OpenML-Referenz vorgibt. Die weitere Bearbeitung soll
EDA-Erkenntnisse ausdrücklich mit den späteren Modellentscheidungen verbinden.}
\field{Schwierigkeiten und offene Fragen}{Die genaue Aufgabenverteilung innerhalb der Gruppe,
die persönliche Forschungsfrage, die Modellstrategie und der verfügbare Rechenrahmen sind noch
zu klären.}
\field{Abweichung vom Vorgehen}{Noch keine Abweichung vom Gruppenplan dokumentiert.}
\field{Gruppentreffen}{\placeholder{kein Treffen / Datum und Kurzverweis ergänzen}}
\field{Zeitaufwand}{\placeholder{Minuten/Stunden ergänzen}}
\field{Nächster Schritt}{Logbuch und E-Portfolio als belastbare Arbeitsgrundlage anlegen und anschließend
den Datensatz laden sowie technisch prüfen.}

\logsection{06.08.2026}{E-Portfolio auf LaTeX umgestellt}
\field{Zusammenfassung}{Das zunächst als Notebook angelegte E-Portfolio wurde auf ein eigenständiges
LaTeX-Dokument umgestellt. Ziel war ein besser kontrollierbares Layout für die spätere PDF-Abgabe.}
\field{Arbeitsschritte}{
\begin{itemize}
  \item Ordner \texttt{E-Protfolio} im eigenen Projektordner angelegt.
  \item Hauptdatei \texttt{ADA\_138\_Sawarzynski\_Portfolio.tex} erstellt.
  \item Titel, Projektüberblick, Logbuch-Platzhalter, Zusammenfassung, Gruppenreflexion,
        Quellen und Abschluss-Checkliste strukturiert.
  \item Ordner \texttt{assets} für spätere Abbildungen vorbereitet.
  \item LaTeX-Datei mit XeLaTeX kompiliert und PDF-Ausgabe geprüft.
\end{itemize}}
\field{Entscheidungen und Begründung}{Das Logbuch wird als separate Datei geführt, damit der
Arbeitsprozess unabhängig vom später verdichteten Portfolio fortlaufend ergänzt werden kann.}
\field{Schwierigkeiten und Lösungen}{Die lokale TinyTeX-Installation enthielt mehrere optionale
Pakete nicht. Die Präambel wurde deshalb auf vorhandene Standardpakete reduziert; die Ausgabe
kompiliert nun erfolgreich.}
\field{Abweichung vom Vorgehen}{Das ursprünglich geplante Notebook-Format wurde zugunsten von
LaTeX verworfen, nachdem die PDF-Gestaltung von nbconvert nicht zufriedenstellend war.}
\field{Gruppentreffen}{Kein inhaltliches Gruppentreffen dokumentiert.}
\field{Zeitaufwand}{\placeholder{Minuten/Stunden ergänzen}}
\field{Nächster Schritt}{Die EDA des Santander-Datensatzes beginnen und dabei Datenform, Datentypen,
Zielvariable, fehlende Werte, Klassenverteilung und erste Verteilungen systematisch untersuchen.}

\logsection{\placeholder{Datum}}{\placeholder{prägnanter Arbeitstitel}}
\field{Zusammenfassung}{\placeholder{Kurze Zusammenfassung des konkreten Arbeitsschritts.}}
\field{Arbeitsschritte}{\placeholder{Welche Schritte wurden tatsächlich durchgeführt?}}
\field{Ergebnis/Erkenntnis}{\placeholder{Was kam dabei heraus? Zahlen, Tabellen oder Abbildungen verlinken.}}
\field{Entscheidung und Begründung}{\placeholder{Welche Entscheidung wurde getroffen und warum?}}
\field{Verworfene Ansätze/Schwierigkeiten}{\placeholder{Was hat nicht funktioniert oder wurde verworfen?}}
\field{Abweichung vom vereinbarten Vorgehen}{\placeholder{Keine / konkrete Abweichung mit Begründung.}}
\field{Gruppentreffen}{\placeholder{Datum, Teilnehmende und Bezug zum Treffen.}}
\field{Zeitaufwand}{\placeholder{Minuten/Stunden.}}
\field{Nächster Schritt}{\placeholder{Konkreter nächster Arbeitsschritt.}}

\section{Laufende Methodik- und Entscheidungsübersicht}

Diese Übersicht wird parallel zum Logbuch aktualisiert. Sie verhindert, dass wichtige
Entscheidungen nur im Fließtext verborgen bleiben.

\begin{longtable}{p{0.19\textwidth}p{0.29\textwidth}p{0.42\textwidth}}
\toprule
\textbf{Bereich} & \textbf{Aktueller Stand} & \textbf{Begründung/Beleg} \\
\midrule
Datenzugang & \placeholder{offen} & \placeholder{Quelle, Version, Zugriffsdatum} \\
Datentypen & \placeholder{offen} & \placeholder{Speicherbedarf und geeignete Wahl} \\
Fehlende Werte & \placeholder{offen} & \placeholder{EDA-Befund und Umgang damit} \\
Zielvariable & \placeholder{offen} & \placeholder{Kodierung und Klassenverteilung} \\
Train-Test-Aufteilung & \placeholder{offen} & \placeholder{Zeitpunkt und Begründung} \\
Metrik(en) & \placeholder{offen} & \placeholder{Warum passend für die Aufgabe?} \\
Modelle & \placeholder{offen} & \placeholder{Vergleich und Auswahl} \\
Feature Engineering & \placeholder{offen} & \placeholder{getestete Varianten} \\
Hyperparameter & \placeholder{offen} & \placeholder{Suchraum und Evaluation} \\
\bottomrule
\end{longtable}

\section{Offene Fragen und Risiken}

\begin{itemize}
  \item \placeholder{Welche genaue persönliche Forschungsfrage wird verfolgt?}
  \item \placeholder{Wie wird die Gruppenarbeit aufgeteilt?}
  \item \placeholder{Welche Metrik ist für die unausgeglichene Zielvariable sachgerecht?}
  \item \placeholder{Wie groß ist der Datensatz und welche Speicherstrategie ist erforderlich?}
  \item \placeholder{Welche Grenzen der Aussagekraft müssen später diskutiert werden?}
\end{itemize}

\section{Hinweise für die spätere Zusammenfassung}

Bei der späteren Verdichtung des Logbuchs müssen Fragestellung, Methodenwahl, Ergebnisse,
Schwierigkeiten/Grenzen und Ausblick klar erkennbar sein. Für Abbildungen und Tabellen werden
Nummer, kurze Caption, Textreferenz und ein prägnanter Take-away ergänzt. Die ausführliche
Nachvollziehbarkeit bleibt in diesem Logbuch erhalten.

\ifdefined\PortfolioMain
  \def\LogbuchEnd{}
\else
  \def\LogbuchEnd{\end{document}}
\fi
\LogbuchEnd


\clearpage
\section{Abschluss-Checkliste}
\begin{itemize}
  \item Name, Gruppennummer, Thema und Gruppenmitglieder ergänzt
  \item Forschungsfrage und Ziel präzisiert
  \item Logbuch chronologisch und zeitnah geführt
  \item Entscheidungen, verworfene Ansätze, Schwierigkeiten und Iterationen dokumentiert
  \item EDA-Erkenntnisse mit Modellentscheidungen verknüpft
  \item Modelle, Metriken, Pipelines, Learning Curves und Hyperparameter beschrieben
  \item Abbildungen/Tabellen nummeriert, beschriftet und interpretiert
  \item Gruppenmitglieder fair und konkret eingeschätzt
  \item Quellen, Softwareversionen und KI-Nutzung vollständig dokumentiert
  \item PDF kompiliert und nach dem vorgegebenen Schema benannt
\end{itemize}

\end{document}
